\documentclass[11pt,a4paper]{article}
\usepackage[utf8]{inputenc}
\usepackage[italian]{babel}
\usepackage{geometry}
\usepackage{graphicx}
\usepackage{hyperref}
\usepackage{listings}
\usepackage{xcolor}
\usepackage{amsmath}
\usepackage{enumitem}
\usepackage{booktabs}
\usepackage{longtable}
\usepackage{array}
\usepackage{tcolorbox}
\usepackage{float}
\usepackage{caption}

\geometry{margin=2.5cm}

% Configurazione codice
\lstset{
    language=JavaScript,
    basicstyle=\ttfamily\small,
    keywordstyle=\color{blue}\bfseries,
    stringstyle=\color{red},
    commentstyle=\color{green!60!black},
    numberstyle=\tiny\color{gray},
    numbers=left,
    numbersep=5pt,
    frame=single,
    breaklines=true,
    showstringspaces=false,
    tabsize=2,
    escapeinside={(*@}{@*)}
}

% Box per evidenziare sezioni importanti
\newtcolorbox{importantbox}[1]{
    colback=red!5!white,
    colframe=red!75!black,
    title=#1,
    fonttitle=\bfseries
}

\newtcolorbox{warningbox}[1]{
    colback=orange!5!white,
    colframe=orange!75!black,
    title=#1,
    fonttitle=\bfseries
}

\newtcolorbox{tipbox}[1]{
    colback=blue!5!white,
    colframe=blue!75!black,
    title=#1,
    fonttitle=\bfseries
}

\newtcolorbox{codebox}[1]{
    colback=gray!5!white,
    colframe=gray!75!black,
    title=#1,
    fonttitle=\bfseries
}

\title{\textbf{Sistema di Studio con Flashcard:\\
Generazione Automatica da PDF}\\
\large Documentazione Tecnica Completa}
\author{Documentazione Sistema ExtraTracker}
\date{\today}

\begin{document}

\maketitle
\tableofcontents
\newpage

\section{Panoramica del Sistema}

\begin{importantbox}{Sistema Critico}
Il sistema di generazione automatica di flashcard da PDF è una \textbf{feature "killer"} che distingue ExtraTracker dalla concorrenza. Utilizza OpenAI GPT per analizzare documenti PDF e generare automaticamente 10-15 flashcard di alta qualità per lo studio.
\end{importantbox}

\subsection{Architettura Generale}

Il sistema è composto da tre layer principali:

\begin{enumerate}
    \item \textbf{Frontend Layer} (React + TypeScript): Componente \texttt{MagicGenerateModal} per l'upload e feedback UX
    \item \textbf{Backend API Layer} (Express.js): Route \texttt{POST /api/study/:id/generate-pdf} con middleware Multer
    \item \textbf{Service Layer} (Node.js): \texttt{StudyService.generateCardsFromPDF()} che orchestrazione estrazione testo, chiamata OpenAI e salvataggio
\end{enumerate}

\subsection{Stack Tecnologico}

\begin{itemize}
    \item \textbf{PDF Parsing}: \texttt{pdf-parse} (npm) per estrazione testo
    \item \textbf{File Upload}: \texttt{multer} per gestione multipart/form-data
    \item \textbf{AI Generation}: OpenAI GPT-4o-mini (configurabile via \texttt{OPENAI\_MODEL})
    \item \textbf{Vector Store}: Pinecone per RAG (Retrieval-Augmented Generation) - opzionale
    \item \textbf{Storage}: MongoDB per persistenza deck e flashcard
    \item \textbf{Frontend PDF Viewer}: \texttt{react-pdf} per visualizzazione interattiva
\end{itemize}

\section{Limitazioni e Vincoli Attuali}

\begin{warningbox}{Limitazioni Attuali}
Il sistema ha attualmente le seguenti limitazioni che devono essere documentate e potenzialmente rimosse in futuro:
\end{warningbox}

\subsection{Limitazione Numero Flashcard}

\textbf{Problema}: Il sistema genera attualmente \textbf{10-15 flashcard} per PDF, indipendentemente dalla lunghezza del documento.

\textbf{Ubicazione nel Codice}:
\begin{lstlisting}[caption=Prompt OpenAI - Limitazione Flashcard]
// server/services/studyService.js:758
const systemPrompt = `...
    REGOLE GENERALI:
    - Crea tra 10 e 15 flashcard
    ...
`;
\end{lstlisting}

\textbf{Motivazione}: 
\begin{itemize}
    \item Controllo costi OpenAI (ogni flashcard generata consuma token)
    \item Qualità vs Quantità: 10-15 flashcard ben fatte sono più utili di 50 flashcard superficiali
    \item Limitazione contesto: \texttt{MAX\_PDF\_TEXT\_LENGTH = 15000} caratteri limita la quantità di contenuto analizzabile
\end{itemize}

\textbf{Impatto}: 
\begin{itemize}
    \item PDF lunghi (100+ pagine) generano solo 10-15 flashcard, coprendo una frazione minima del contenuto
    \item Utenti devono caricare lo stesso PDF più volte o dividere manualmente il documento
\end{itemize}

\subsection{Limitazione Lunghezza Testo}

\textbf{Problema}: Il testo estratto dal PDF viene troncato a \textbf{15.000 caratteri} prima dell'invio a OpenAI.

\textbf{Ubicazione nel Codice}:
\begin{lstlisting}[caption=Costante e Troncamento Testo]
// server/services/studyService.js:27
const MAX_PDF_TEXT_LENGTH = 15000; // Limite caratteri per evitare costi eccessivi

// server/services/studyService.js:717-719
const truncatedText = normalizedExtracted.length > MAX_PDF_TEXT_LENGTH 
    ? normalizedExtracted.slice(0, MAX_PDF_TEXT_LENGTH) + '\n\n[...testo troncato...]'
    : normalizedExtracted;
\end{lstlisting}

\textbf{Motivazione}:
\begin{itemize}
    \item Controllo costi: più testo = più token = più costi OpenAI
    \item Limiti modello: GPT-4o-mini ha context window di 128k token, ma 15k caratteri è un limite conservativo
    \item Performance: testi più lunghi richiedono più tempo di elaborazione
\end{itemize}

\textbf{Impatto}:
\begin{itemize}
    \item PDF con più di ~10-15 pagine (dipende dalla densità di testo) vengono troncati
    \item Le flashcard generate coprono solo le prime pagine del documento
    \item Informazioni nelle pagine successive vengono ignorate
\end{itemize}

\subsection{Limitazione Dimensione File}

\textbf{Problema}: Il sistema accetta PDF fino a \textbf{10MB}.

\textbf{Ubicazione nel Codice}:
\begin{lstlisting}[caption=Validazione Dimensione File - Frontend]
// src/features/study/services/studyService.ts:368-370
if (file.size > 10 * 1024 * 1024) {
    throw new Error('Il file supera il limite di 10MB');
}
\end{lstlisting}

\begin{lstlisting}[caption=Validazione Dimensione File - Backend]
// server/routes/study.js:46-48
const upload = multer({
    storage,
    limits: {
        fileSize: 10 * 1024 * 1024, // 10MB max
    },
    ...
});
\end{lstlisting}

\textbf{Motivazione}:
\begin{itemize}
    \item Storage: limitare lo spazio disco utilizzato
    \item Performance: file più grandi richiedono più tempo per parsing e upload
    \item Network: upload di file grandi può fallire su connessioni lente
\end{itemize}

\section{Flusso Completo: Upload PDF → Flashcard}

\subsection{Diagramma di Flusso}

Il processo completo può essere visualizzato come segue:

\begin{enumerate}
    \item \textbf{Upload Frontend} → Utente seleziona/droppa PDF in \texttt{MagicGenerateModal}
    \item \textbf{Validazione Client-Side} → Verifica tipo file (PDF), dimensione (≤10MB)
    \item \textbf{FormData Upload} → \texttt{POST /api/study/:id/generate-pdf} con \texttt{multipart/form-data}
    \item \textbf{Middleware Multer} → Salvataggio file su disco in \texttt{server/uploads/pdfs/}
    \item \textbf{Controller Validation} → Verifica file ricevuto, tipo MIME, dimensione
    \item \textbf{Service Processing} → \texttt{generateCardsFromPDF()} esegue:
        \begin{itemize}
            \item Lettura file PDF da disco
            \item Estrazione testo con \texttt{pdf-parse}
            \item Normalizzazione e formattazione testo (con marker pagine)
            \item Troncamento a 15k caratteri se necessario
            \item Chiamata OpenAI con prompt personalizzato
            \item Parsing risposta JSON
            \item Validazione e normalizzazione flashcard
            \item Salvataggio bulk nel deck MongoDB
        \end{itemize}
    \item \textbf{Response} → Ritorna deck aggiornato con conteggio flashcard generate
    \item \textbf{Frontend Update} → UI aggiornata con nuove flashcard
\end{enumerate}

\subsection{Componente Frontend: MagicGenerateModal}

\subsubsection{Responsabilità}

Il componente \texttt{MagicGenerateModal.tsx} gestisce:
\begin{itemize}
    \item \textbf{UI/UX}: Drag \& drop area, progress indicator multi-step, feedback visivo
    \item \textbf{Validazione Client-Side}: Tipo file, dimensione
    \item \textbf{State Management}: File selezionato, step corrente (idle → uploading → reading → analyzing → generating → success)
    \item \textbf{Error Handling}: Gestione errori con toast e messaggi user-friendly
\end{itemize}

\subsubsection{Stati del Processo}

\begin{lstlisting}[caption=Stati Generazione - Frontend]
type GenerationStep = 
    | 'idle'           // Nessun file selezionato
    | 'uploading'      // Upload in corso
    | 'reading'        // Lettura PDF sul server
    | 'analyzing'      // Analisi contenuti
    | 'generating'     // Generazione flashcard con AI
    | 'success'        // Completato con successo
    | 'error';         // Errore durante il processo
\end{lstlisting}

\subsubsection{Simulazione Progress}

\begin{importantbox}{Nota Implementazione}
Il frontend \textbf{simula} i passaggi di progress per migliorare l'UX, ma i tempi reali dipendono dal backend. La simulazione avviene in parallelo con la chiamata API reale.
\end{importantbox}

\begin{lstlisting}[caption=Simulazione Progress Steps]
// src/features/study/components/MagicGenerateModal.tsx:125-136
const simulateSteps = async () => {
    setStep('uploading');
    await new Promise(r => setTimeout(r, 800));
    
    setStep('reading');
    await new Promise(r => setTimeout(r, 1200));
    
    setStep('analyzing');
    await new Promise(r => setTimeout(r, 1500));
    
    setStep('generating');
};
\end{lstlisting}

\subsubsection{Chiamata API}

\begin{lstlisting}[caption=Upload PDF - Frontend Service]
// src/features/study/services/studyService.ts:360-410
async generateFromPDF(deckId: string, file: File) {
    // Validazione client-side
    if (file.type !== 'application/pdf') {
        throw new Error('Solo file PDF sono supportati');
    }
    if (file.size > 10 * 1024 * 1024) {
        throw new Error('Il file supera il limite di 10MB');
    }

    // Costruisci FormData
    const formData = new FormData();
    formData.append('pdf', file);

    // Chiamata API con credentials per cookies HttpOnly
    const response = await fetch(`/api/study/${deckId}/generate-pdf`, {
        method: 'POST',
        body: formData,
        credentials: 'include', // Include cookies per auth
    });

    if (!response.ok) {
        const errorData = await response.json().catch(() => ({}));
        throw new Error(errorData.error?.message || 'Generazione fallita');
    }

    const result = await response.json();
    return {
        generatedCount: result.data?.generatedCount || 0,
        deck: normalizeDeck(result.data?.deck || {}),
    };
}
\end{lstlisting}

\subsection{Backend Route: Upload Handler}

\subsubsection{Multer Configuration}

\begin{lstlisting}[caption=Configurazione Multer per PDF Upload]
// server/routes/study.js:25-56
const PDF_UPLOAD_DIR = path.join(__dirname, '..', 'uploads', 'pdfs');
fs.mkdirSync(PDF_UPLOAD_DIR, { recursive: true });

const storage = multer.diskStorage({
    destination: (req, _file, cb) => {
        fs.mkdirSync(PDF_UPLOAD_DIR, { recursive: true });
        cb(null, PDF_UPLOAD_DIR);
    },
    filename: (req, _file, cb) => {
        const deckId = req.params?.id || 'deck';
        const timestamp = Date.now();
        cb(null, `${deckId}-${timestamp}.pdf`);
    },
});

const upload = multer({
    storage,
    limits: {
        fileSize: 10 * 1024 * 1024, // 10MB max
    },
    fileFilter: (req, file, cb) => {
        if (file.mimetype === 'application/pdf') {
            cb(null, true);
        } else {
            cb(new Error('Solo file PDF sono accettati'), false);
        }
    },
});
\end{lstlisting}

\begin{tipbox}{Naming Convention File}
I file PDF vengono salvati con il pattern: \texttt{\{deckId\}-\{timestamp\}.pdf}
Esempio: \texttt{507f1f77bcf86cd799439011-1767551908605.pdf}

Questo garantisce:
\begin{itemize}
    \item Unicità: timestamp evita collisioni
    \item Tracciabilità: deckId permette di associare file a deck
    \item Persistenza: file rimangono sul server anche se la generazione fallisce
\end{itemize}
\end{tipbox}

\subsubsection{Route Handler}

\begin{lstlisting}[caption=Route Handler con Multer Middleware]
// server/routes/study.js:84-97
router.post('/:id/generate-pdf', (req, res, next) => {
    upload.single('pdf')(req, res, (err) => {
        if (err) {
            console.error('Multer error:', err.message);
            return res.status(400).json({
                success: false,
                error: { message: err.message || 'Errore nel caricamento del file' }
            });
        }
        console.log('📄 File ricevuto:', req.file ? req.file.originalname : 'NESSUN FILE');
        next();
    });
}, studyController.uploadAndGenerate);
\end{lstlisting}

\subsubsection{Controller Validation}

\begin{lstlisting}[caption=Controller - Validazione e Chiamata Service]
// server/controllers/studyController.js:189-232
const uploadAndGenerate = asyncHandler(async (req, res) => {
    // Verifica che il file sia stato caricato
    if (!req.file) {
        return res.status(400).json({
            success: false,
            error: { message: 'Nessun file caricato. Carica un file PDF.' },
        });
    }

    // Verifica che sia un PDF
    if (req.file.mimetype !== 'application/pdf') {
        return res.status(400).json({
            success: false,
            error: 'Il file deve essere un PDF.',
        });
    }

    // Limite dimensione file (10MB)
    const maxSize = 10 * 1024 * 1024;
    if (req.file.size > maxSize) {
        return res.status(400).json({
            success: false,
            error: 'Il file è troppo grande. Massimo 10MB.',
        });
    }

    // Chiamata service per generazione
    const result = await studyService.generateCardsFromPDF(
        req.tenantScope,
        req.params.id,
        req.file.path
    );

    res.json({
        success: true,
        data: result,
        message: `✨ Generate ${result.generatedCount} flashcard con successo!`,
    });
});
\end{lstlisting}

\section{Service Layer: Estrazione e Generazione}

\subsection{Metodo Principale: generateCardsFromPDF}

\begin{lstlisting}[caption=Metodo Principale - generateCardsFromPDF]
// server/services/studyService.js:659-854
async generateCardsFromPDF(tenantScope, deckId, pdfFilePath) {
    const userId = this._getUserId(tenantScope);

    // 1. Verifica ownership deck
    const deck = await Deck.findOne({ _id: deckId, user: userId });
    if (!deck) {
        throw AppError.notFound('Mazzo');
    }

    // 2. Leggi file PDF da disco
    let pdfBuffer;
    try {
        pdfBuffer = await fs.readFile(pdfFilePath);
    } catch (err) {
        throw AppError.validation('Impossibile leggere il PDF caricato.');
    }

    // 3. Estrai testo dal PDF
    let pdfText;
    try {
        const pdfData = await pdfParse(pdfBuffer);
        pdfText = this._formatPdfTextWithPages(pdfData);
    } catch (err) {
        throw AppError.validation('Impossibile leggere il PDF. Assicurati che sia valido e non protetto.');
    }

    if (!pdfText || pdfText.trim().length < 50) {
        throw AppError.validation('Il PDF non contiene abbastanza testo da elaborare.');
    }

    // 4. Normalizza e salva testo estratto
    const normalizedExtracted = this._normalizeExtractedText(pdfText);
    deck.extractedText = this._truncateText(
        normalizedExtracted,
        MAX_EXTRACTED_TEXT_STORE_LENGTH, // 200k caratteri
        '\n\n[...testo troncato per limiti di storage...]'
    );
    await deck.save({ validateModifiedOnly: true });

    // 5. Ingestione vettoriale (best-effort, non blocca)
    try {
        await vectorStoreService.ingestDeck(deckId, normalizedExtracted);
    } catch (err) {
        console.error('⚠️ Vector ingest error:', err.message);
    }

    // 6. Tronca testo per OpenAI (limite 15k caratteri)
    const truncatedText = normalizedExtracted.length > MAX_PDF_TEXT_LENGTH 
        ? normalizedExtracted.slice(0, MAX_PDF_TEXT_LENGTH) + '\n\n[...testo troncato...]'
        : normalizedExtracted;

    // 7. Genera flashcard con OpenAI
    const aiResponse = await this._callOpenAIForCards(truncatedText, deck);

    // 8. Parsa e valida flashcard
    const validCards = this._parseAndValidateCards(aiResponse);

    // 9. Salva flashcard nel deck
    deck.cards.push(...validCards);
    await deck.save();

    return {
        deck: deck.toJSON(),
        generatedCount: validCards.length,
    };
}
\end{lstlisting}

\subsection{Estrazione Testo PDF}

\subsubsection{Libreria pdf-parse}

Il sistema utilizza \texttt{pdf-parse} per estrarre testo da PDF. Questa libreria:
\begin{itemize}
    \item Supporta PDF con testo selezionabile (non funziona con PDF scansionati/immagini)
    \item Estrae testo pagina per pagina
    \item Mantiene struttura base (paragrafi, linee)
    \item Non preserva formattazione complessa (tabelle, colonne multiple)
\end{itemize}

\subsubsection{Formattazione con Marker Pagine}

\begin{lstlisting}[caption=Formattazione Testo con Marker Pagine]
// server/services/studyService.js:1340-1353
_formatPdfTextWithPages(pdfTextResult) {
    const pages = Array.isArray(pdfTextResult?.pages) ? pdfTextResult.pages : [];
    const fallback = typeof pdfTextResult?.text === 'string' ? pdfTextResult.text : '';

    if (pages.length === 0) return fallback;

    return pages
        .map((page, index) => {
            const pageNumber = Number.isFinite(Number(page?.num)) 
                ? Number(page.num) 
                : index + 1;
            const text = typeof page?.text === 'string' ? page.text : '';
            return `\n--- Pagina ${pageNumber} ---\n${text}`;
        })
        .join('\n');
}
\end{lstlisting}

\begin{tipbox}{Marker Pagine per RAG}
I marker di pagina (\texttt{--- Pagina X ---}) sono cruciali per:
\begin{itemize}
    \item \textbf{AI Tutor}: Permette all'AI di rispondere a domande tipo "spiegami pagina 5"
    \item \textbf{Tracciabilità}: Utente può sapere da quale pagina proviene una flashcard
    \item \textbf{Debug}: Facilita identificazione problemi di estrazione
\end{itemize}
\end{tipbox}

\subsubsection{Normalizzazione Testo}

\begin{lstlisting}[caption=Normalizzazione Testo Estratto]
// server/services/studyService.js:1332-1338
_normalizeExtractedText(text) {
    if (!text || typeof text !== 'string') return '';
    return text
        .replace(/\r\n/g, '\n')           // Normalizza line breaks
        .replace(/\n{3,}/g, '\n\n')       // Rimuove spazi multipli
        .trim();
}
\end{lstlisting}

\subsection{Chiamata OpenAI per Generazione}

\subsubsection{Prompt Engineering}

Il sistema utilizza un prompt strutturato che include:
\begin{itemize}
    \item \textbf{System Prompt}: Definisce il ruolo dell'AI (esperto insegnante universitario "Silvija")
    \item \textbf{Style Prompts}: Configurazione stile (comprehensive, conceptual, factual, application)
    \item \textbf{Difficulty Prompts}: Configurazione difficoltà (easy, medium, hard, mixed)
    \item \textbf{Question Type Prompts}: Tipi di domande da generare (definition, concept, relationship, etc.)
    \item \textbf{Regole Generali}: Istruzioni per qualità flashcard
    \item \textbf{Formato Output}: JSON strutturato con array di cards
\end{itemize}

\begin{lstlisting}[caption=Costruzione System Prompt]
// server/services/studyService.js:751-767
const systemPrompt = `Sei un esperto insegnante universitario che si chiama Silvija. 
Il tuo compito è analizzare il testo fornito e creare flashcard di alta qualità.

STILE: ${stylePrompts[style] || stylePrompts.comprehensive}
DIFFICOLTÀ: ${difficultyPrompts[difficulty] || difficultyPrompts.medium}
TIPI DI DOMANDE: ${questionTypes.map(t => questionTypePrompts[t]).join(' ')}

REGOLE GENERALI:
- Crea tra 10 e 15 flashcard
- Ogni flashcard deve avere "front" (domanda) e "back" (risposta)
- Le domande devono testare comprensione, non solo memoria
- Evita domande troppo generiche o troppo specifiche
- La risposta deve essere auto-contenuta

FORMATO OUTPUT:
Rispondi SOLO con JSON valido, senza markdown:
{"cards":[{"front":"Domanda 1?","back":"Risposta 1"},...]}`;
\end{lstlisting}

\subsubsection{Chiamata API OpenAI}

\begin{lstlisting}[caption=Chiamata OpenAI Chat Completions]
// server/services/studyService.js:769-797
const completion = await openai.chat.completions.create({
    model: process.env.OPENAI_MODEL || 'gpt-4o-mini',
    messages: [
        { role: 'system', content: systemPrompt },
        { role: 'user', content: `Analizza questo testo e genera le flashcard:\n\n${truncatedText}` }
    ],
    temperature: 0.7,                    // Bilanciamento creatività/precisione
    max_completion_tokens: 2000,         // Limite token risposta
    response_format: { type: 'json_object' }, // Forza output JSON
});
\end{lstlisting}

\begin{importantbox}{Configurazione Modello}
Il modello OpenAI è configurabile via variabile ambiente \texttt{OPENAI\_MODEL}. Default: \texttt{gpt-4o-mini}

Modelli supportati:
\begin{itemize}
    \item \texttt{gpt-4o-mini}: Economico, veloce, buona qualità (consigliato)
    \item \texttt{gpt-4o}: Più costoso, migliore qualità, più lento
    \item \texttt{gpt-3.5-turbo}: Legacy, non consigliato
\end{itemize}
\end{importantbox}

\subsubsection{Parsing e Validazione Risposta}

\begin{lstlisting}[caption=Parsing Risposta OpenAI]
// server/services/studyService.js:799-842
// 1. Parsa JSON
let generatedCards;
try {
    const parsed = JSON.parse(aiResponse);
    generatedCards = this._extractGeneratedCards(parsed);
} catch (err) {
    throw AppError.internal('Risposta AI non valida. Riprova.');
}

// 2. Estrai cards da vari formati possibili
_extractGeneratedCards(parsed) {
    if (Array.isArray(parsed)) return parsed;
    if (Array.isArray(parsed.cards)) return parsed.cards;
    if (Array.isArray(parsed.flashcards)) return parsed.flashcards;
    if (Array.isArray(parsed.items)) return parsed.items;
    if (parsed.front && parsed.back) return [parsed];
    // ... altri fallback
}

// 3. Normalizza formato card
_normalizeGeneratedCard(card) {
    return {
        front: card.front ?? card.question ?? card.q,
        back: card.back ?? card.answer ?? card.a,
    };
}

// 4. Valida e filtra cards valide
const validCards = generatedCards
    .map((card) => this._normalizeGeneratedCard(card))
    .filter(card => card.front && card.back && 
                    typeof card.front === 'string' && 
                    typeof card.back === 'string')
    .map(card => ({
        front: card.front.trim(),
        back: card.back.trim(),
        status: 'new',
        nextReviewDate: new Date(),
        easinessFactor: DEFAULT_EASINESS_FACTOR, // 2.5
        interval: 0,
        repetitions: 0,
    }));
\end{lstlisting}

\section{Vector Store Integration (RAG)}

\subsection{Ingestione Vettoriale}

Dopo l'estrazione del testo, il sistema tenta (best-effort) di ingerire il testo in Pinecone per abilitare RAG:

\begin{lstlisting}[caption=Ingestione Vettoriale]
// server/services/studyService.js:709-714
try {
    await vectorStoreService.ingestDeck(deckId, normalizedExtracted);
} catch (err) {
    console.error('⚠️ Vector ingest error:', err.message);
    // Non blocca il processo se fallisce
}
\end{lstlisting}

\begin{tipbox}{Best-Effort Pattern}
L'ingestione vettoriale è \textbf{best-effort}: se fallisce, il processo continua normalmente. Questo perché:
\begin{itemize}
    \item RAG è opzionale (usato solo per AI Tutor)
    \item Pinecone potrebbe non essere configurato
    \item Errori di rete non devono bloccare la generazione flashcard
\end{itemize}
\end{tipbox}

\subsection{Configurazione Vector Store}

\begin{lstlisting}[caption=Configurazione Pinecone]
// server/services/vectorStoreService.js:13-14
const DEFAULT_EMBED_MODEL = process.env.OPENAI_EMBED_MODEL || 'text-embedding-3-small';
const DEFAULT_TOP_K = 5;

// Chunking: RecursiveCharacterTextSplitter
// - Chunk size: 1000 caratteri
// - Overlap: 200 caratteri
// - Embedding: OpenAI text-embedding-3-small
// - Vector DB: Pinecone (filtrato per deckId)
\end{lstlisting}

\section{Modello Dati: Deck e Cards}

\subsection{Schema Deck}

\begin{lstlisting}[caption=Schema Deck - Campi Rilevanti PDF]
// server/models/Deck.js:110-119
const deckSchema = new mongoose.Schema({
    // ...
    pdfUrl: {
        type: String,
        default: '',
        trim: true,
    },
    extractedText: {
        type: String,
        default: '',
        select: false,  // Non incluso di default nelle query
    },
    // ...
});
\end{lstlisting}

\begin{importantbox}{Campo extractedText}
Il campo \texttt{extractedText} è marcato con \texttt{select: false} per:
\begin{itemize}
    \item \textbf{Performance}: Testi estratti possono essere molto grandi (fino a 200k caratteri)
    \item \textbf{Bandwidth}: Evita di inviare dati non necessari al frontend
    \item \textbf{Privacy}: Testi completi non vengono esposti nelle API standard
\end{itemize}

Per accedere al testo estratto, usa query esplicita: \texttt{Deck.findById(id).select('+extractedText')}
\end{importantbox}

\subsection{Schema Card}

\begin{lstlisting}[caption=Schema Card - Flashcard Generata]
// server/models/Deck.js:19-81
const cardSchema = new mongoose.Schema({
    front: {
        type: String,
        required: true,
        trim: true,
        maxlength: [2000, 'Il fronte non può superare 2000 caratteri'],
    },
    back: {
        type: String,
        required: true,
        trim: true,
        maxlength: [4000, 'Il retro non può superare 4000 caratteri'],
    },
    // SM-2 parameters (per spaced repetition)
    easinessFactor: { type: Number, default: 2.5, min: 1.3 },
    interval: { type: Number, default: 0, min: 0 },
    repetitions: { type: Number, default: 0, min: 0 },
    nextReviewDate: { type: Date, default: Date.now },
    status: {
        type: String,
        enum: ['new', 'learning', 'review', 'mastered'],
        default: 'new',
    },
    // ...
});
\end{lstlisting}

\section{Error Handling e Edge Cases}

\subsection{Errori Comuni}

\begin{enumerate}
    \item \textbf{PDF Protetto/Password}: \texttt{pdf-parse} fallisce → Errore validazione
    \item \textbf{PDF Scansionato (Solo Immagini)}: Nessun testo estratto → Errore "non contiene abbastanza testo"
    \item \textbf{PDF Vuoto}: Testo < 50 caratteri → Errore validazione
    \item \textbf{OpenAI Quota Esaurita}: \texttt{insufficient\_quota} → Errore interno con messaggio user-friendly
    \item \textbf{OpenAI Context Length}: Testo troppo lungo → Errore (raro, testo già troncato a 15k)
    \item \textbf{JSON Parsing Error}: Risposta AI non valida → Errore interno
    \item \textbf{File Upload Timeout}: Upload > 60s → Errore network
\end{enumerate}

\subsection{Gestione Errori nel Frontend}

\begin{lstlisting}[caption=Error Handling - Frontend]
// src/features/study/components/MagicGenerateModal.tsx:158-162
catch (err: any) {
    setStep('error');
    setError(err.message || 'Errore nella generazione. Riprova.');
    emitToast.error(err.message || 'Generazione fallita');
}
\end{lstlisting}

\subsection{Gestione Errori nel Backend}

\begin{lstlisting}[caption=Error Handling - Backend]
// Tutti gli errori vengono wrappati in AppError
// server/utils/AppError.js gestisce formattazione consistente

// Esempi:
throw AppError.validation('Il PDF non contiene abbastanza testo.');
throw AppError.internal('Errore nella generazione AI.', err, {});
throw AppError.notFound('Mazzo');
\end{lstlisting}

\section{Configurazione AI Settings}

\subsection{Impostazioni Personalizzabili}

Il sistema permette di configurare la generazione tramite \texttt{deck.aiSettings}:

\begin{lstlisting}[caption=AI Settings Schema]
// server/models/Deck.js:137-152
aiSettings: {
    style: {
        type: String,
        enum: ['comprehensive', 'conceptual', 'factual', 'application'],
        default: 'comprehensive',
    },
    difficulty: {
        type: String,
        enum: ['easy', 'medium', 'hard', 'mixed'],
        default: 'medium',
    },
    questionTypes: {
        type: [String],
        default: ['definition', 'concept', 'relationship'],
    },
}
\end{lstlisting}

\subsection{Style Prompts}

\begin{itemize}
    \item \textbf{comprehensive}: Bilanciato tra teoria e pratica
    \item \textbf{conceptual}: Focus su concetti e principi
    \item \textbf{factual}: Focus su fatti, date, definizioni
    \item \textbf{application}: Focus su applicazioni pratiche
\end{itemize}

\subsection{Difficulty Prompts}

\begin{itemize}
    \item \textbf{easy}: Domande semplici per principianti
    \item \textbf{medium}: Difficoltà media (default)
    \item \textbf{hard}: Domande complesse che richiedono analisi
    \item \textbf{mixed}: Mix di difficoltà
\end{itemize}

\subsection{Question Types}

\begin{itemize}
    \item \textbf{definition}: Definizioni precise
    \item \textbf{concept}: Comprensione concetti astratti
    \item \textbf{relationship}: Relazioni, cause-effetti
    \item \textbf{application}: Applicazioni pratiche
    \item \textbf{comparison}: Confronti e distinzioni
\end{itemize}

\section{Performance e Ottimizzazioni}

\subsection{Tempi Stimati}

\begin{itemize}
    \item \textbf{Upload File}: 1-5s (dipende da dimensione e connessione)
    \item \textbf{Parsing PDF}: 0.5-2s (dipende da numero pagine)
    \item \textbf{Estrazione Testo}: 0.1-1s
    \item \textbf{Chiamata OpenAI}: 3-10s (dipende da modello e lunghezza testo)
    \item \textbf{Salvataggio DB}: 0.1-0.5s
    \item \textbf{Totale}: 5-20s tipicamente
\end{itemize}

\subsection{Bottlenecks Potenziali}

\begin{enumerate}
    \item \textbf{OpenAI API}: Rate limiting e latenza possono rallentare
    \item \textbf{Parsing PDF Grandi}: PDF con 100+ pagine richiedono più tempo
    \item \textbf{Vector Ingest}: Pinecone può essere lento per testi lunghi
    \item \textbf{DB Write}: Salvataggio bulk di molte flashcard può essere lento
\end{enumerate}

\subsection{Ottimizzazioni Implementate}

\begin{itemize}
    \item \textbf{Troncamento Testo}: Limite 15k caratteri riduce costi e latenza OpenAI
    \item \textbf{Bulk Insert}: Flashcard salvate in un'unica operazione
    \item \textbf{Best-Effort Vector}: Ingestione vettoriale non blocca il processo
    \item \textbf{Async Processing}: Frontend non blocca durante generazione
\end{itemize}

\section{Possibili Miglioramenti Futuri}

\subsection{Rimozione Limitazione 10-15 Flashcard}

\textbf{Problema Attuale}: Solo 10-15 flashcard generate indipendentemente dalla lunghezza PDF.

\textbf{Soluzione Proposta}:
\begin{enumerate}
    \item \textbf{Calcolo Dinamico}: Numero flashcard basato su lunghezza testo
        \begin{itemize}
            \item 1 flashcard ogni 500-1000 caratteri
            \item Minimo 10, massimo 50-100
        \end{itemize}
    \item \textbf{Chunking Intelligente}: Dividere PDF in sezioni e generare flashcard per sezione
    \item \textbf{Multi-Pass Generation}: 
        \begin{itemize}
            \item Pass 1: Analisi struttura documento
            \item Pass 2: Generazione flashcard per ogni sezione importante
        \end{itemize}
    \item \textbf{User Control}: Permettere all'utente di specificare numero desiderato
\end{enumerate}

\subsection{Rimozione Limitazione 15k Caratteri}

\textbf{Problema Attuale}: Solo primi 15k caratteri analizzati.

\textbf{Soluzione Proposta}:
\begin{enumerate}
    \item \textbf{Chunking Strategico}: 
        \begin{itemize}
            \item Dividere PDF in chunk di 10-15k caratteri
            \item Generare flashcard per ogni chunk
            \item Unire risultati
        \end{itemize}
    \item \textbf{Selezione Intelligente}:
        \begin{itemize}
            \item Analisi preliminare per identificare sezioni più importanti
            \item Generare flashcard solo per sezioni rilevanti
        \end{itemize}
    \item \textbf{Summary + Details}:
        \begin{itemize}
            \item Generare flashcard di riepilogo per tutto il documento
            \item Generare flashcard dettagliate per sezioni specifiche
        \end{itemize}
\end{enumerate}

\subsection{Supporto PDF Scansionati (OCR)}

\textbf{Problema Attuale}: PDF scansionati (immagini) non funzionano.

\textbf{Soluzione Proposta}:
\begin{enumerate}
    \item \textbf{Integrazione OCR}: Usare Tesseract.js o servizio cloud (Google Vision, AWS Textract)
    \item \textbf{Pre-Processing}: Rilevare se PDF contiene solo immagini
    \item \textbf{Fallback Automatico}: Se estrazione testo fallisce, tentare OCR
\end{enumerate}

\subsection{Progress Tracking Reale}

\textbf{Problema Attuale}: Frontend simula progress, non riflette stato reale backend.

\textbf{Soluzione Proposta}:
\begin{enumerate}
    \item \textbf{WebSocket/SSE}: Comunicazione real-time backend → frontend
    \item \textbf{Progress Events}: Eventi per ogni step (parsing, extraction, generation)
    \item \textbf{Percentuale Reale}: Calcolo percentuale basato su operazioni completate
\end{enumerate}

\subsection{Preview Flashcard Prima di Salvare}

\textbf{Problema Attuale}: Utente non può vedere flashcard prima di salvarle.

\textbf{Soluzione Proposta}:
\begin{enumerate}
    \item \textbf{Modal Preview}: Mostrare flashcard generate in modal
    \item \textbf{Edit Before Save}: Permettere modifica/eliminazione flashcard
    \item \textbf{Selective Save}: Salvare solo flashcard selezionate
\end{enumerate}

\section{Testing e Debug}

\subsection{Logging}

Il sistema include logging estensivo per debug:

\begin{lstlisting}[caption=Logging Points]
// Upload
console.log('📤 Uploading PDF:', file.name, file.size, 'bytes');
console.log('📥 Response status:', response.status);

// Backend
console.log('🪄 uploadAndGenerate called');
console.log('📄 File ricevuto:', req.file?.originalname);
console.log('✨ Generated ${validCards.length} flashcards');

// Errori
console.error('❌ PDF Parse Error:', err.message);
console.error('❌ OpenAI API Error:', err.message);
console.error('⚠️ Vector ingest error:', err.message);
\end{lstlisting}

\subsection{Test Manuali}

Per testare il sistema:

\begin{enumerate}
    \item \textbf{PDF Valido}: Caricare PDF con testo selezionabile (10-20 pagine)
    \item \textbf{PDF Protetto}: Tentare caricamento PDF con password → Verificare errore
    \item \textbf{PDF Grande}: Caricare PDF > 10MB → Verificare errore
    \item \textbf{PDF Scansionato}: Caricare PDF solo immagini → Verificare errore "non contiene testo"
    \item \textbf{PDF Vuoto}: Caricare PDF vuoto → Verificare errore
\end{enumerate}

\section{Conclusioni}

Il sistema di generazione automatica di flashcard da PDF è un componente complesso che integra:
\begin{itemize}
    \item Upload e storage file
    \item Estrazione testo da PDF
    \item Chiamate API OpenAI
    \item Parsing e validazione risposte
    \item Persistenza dati MongoDB
    \item Integrazione vettoriale (RAG)
\end{itemize}

Le limitazioni attuali (10-15 flashcard, 15k caratteri) sono progettate per bilanciare:
\begin{itemize}
    \item \textbf{Costi}: Controllo spese OpenAI
    \item \textbf{Performance}: Tempi di risposta accettabili
    \item \textbf{Qualità}: Flashcard ben fatte vs quantità
\end{itemize}

Future iterazioni possono rimuovere queste limitazioni implementando chunking intelligente, multi-pass generation, e controlli utente più granulari.

\end{document}
